\begingroup
\let\cleardoublepage\relax


\chapter*{Operatoren in Java}
\label{ch:operatoren-in-java}
\addcontentsline{toc}{chapter}{Operatoren in Java}
\endgroup


\section*{Arithmetische Operatoren}
\addcontentsline{toc}{section}{Arithmetische Operatoren}

\subsection*{Beschreibung:}\\
Werden für mathematische Berechnungen verwendet.

\begin{center}
\begin{tabular}{|l|l|l|l|}
\hline
\textbf{Operator} & \textbf{Bedeutung} & \textbf{Beispiel} & \textbf{Ergebnis} \\
\hline
\texttt{+} & Addition & \texttt{5 + 3} & \texttt{8} \\
\hline
\texttt{-} & Subtraktion & \texttt{5 - 3} & \texttt{2} \\
\hline
\texttt{*} & Multiplikation & \texttt{5 * 3} & \texttt{15} \\
\hline
\texttt{/} & Division & \texttt{10 / 3} & \texttt{3} (bei int) \\
\hline
\texttt{\%} & Modulo (Rest) & \texttt{10 \% 3} & \texttt{1} \\
\hline
\end{tabular}
\end{center}

\textbf{Beispiele:}
\begin{lstlisting}
int a = 10;
int b = 3;
int summe = a + b;        // summe = 13
int differenz = a - b;    // differenz = 7
int produkt = a * b;      // produkt = 30
int quotient = a / b;     // quotient = 3 (Ganzzahldivision!)
int rest = a % b;         // rest = 1

// Gerade oder ungerade pruefen
int zahl = 7;
if (zahl % 2 == 0) {
    System.out.println("Gerade");
} else {
    System.out.println("Ungerade");
}
\end{lstlisting}

\textbf{Wichtig bei Division:}
\begin{lstlisting}
int x = 7 / 2;      // x = 3 (int-Division schneidet ab)
double y = 7.0 / 2; // y = 3.5 (mit Kommazahl)
\end{lstlisting}

% ============================================================
\newpage


\section*{Vergleichsoperatoren}
\addcontentsline{toc}{section}{Vergleichsoperatoren}

\subsection*{Beschreibung:}\\
Vergleichen zwei Werte und geben \texttt{true} oder \texttt{false} zurück.

\begin{center}
\begin{tabular}{|l|l|l|l|}
\hline
\textbf{Operator} & \textbf{Bedeutung} & \textbf{Beispiel} & \textbf{Ergebnis} \\
\hline
\texttt{==} & Gleich & \texttt{5 == 5} & \texttt{true} \\
\hline
\texttt{!=} & Ungleich & \texttt{5 != 3} & \texttt{true} \\
\hline
\texttt{>} & Größer & \texttt{5 > 3} & \texttt{true} \\
\hline
\texttt{<} & Kleiner & \texttt{5 < 3} & \texttt{false} \\
\hline
\texttt{>=} & Größer oder gleich & \texttt{5 >= 5} & \texttt{true} \\
\hline
\texttt{<=} & Kleiner oder gleich & \texttt{5 <= 3} & \texttt{false} \\
\hline
\end{tabular}
\end{center}

\textbf{Beispiele:}
\begin{lstlisting}
int a = 10;
int b = 5;

if (a > b) {
    System.out.println("a ist groesser");
}

if (a == 10) {
    System.out.println("a ist 10");
}

if (b != 0) {
    int ergebnis = a / b;  // Division nur wenn b nicht 0 ist
}

// Array-Index pruefen
int[] zahlen = {1, 2, 3};
int index = 2;
if (index >= 0 && index < zahlen.length) {
    System.out.println(zahlen[index]);
}
\end{lstlisting}

\textbf{Wichtig bei Strings:}
\begin{lstlisting}
String text1 = "Hallo";
String text2 = "Hallo";

// FALSCH - vergleicht Objektreferenzen, nicht Inhalt!
if (text1 == text2) { }

// RICHTIG - vergleicht Inhalt
if (text1.equals(text2)) { }
\end{lstlisting}

% ============================================================
\newpage


\section*{Logische Operatoren}
\addcontentsline{toc}{section}{Logische Operatoren}

\subsection*{Beschreibung:}\\
Verknüpfen mehrere Bedingungen.

\begin{center}
\begin{tabular}{|l|l|l|l|}
\hline
\textbf{Operator} & \textbf{Bedeutung} & \textbf{Beispiel} & \textbf{Ergebnis} \\
\hline
\texttt{\&\&} & Und (AND) & \texttt{true \&\& false} & \texttt{false} \\
\hline
\texttt{||} & Oder (OR) & \texttt{true || false} & \texttt{true} \\
\hline
\texttt{!} & Nicht (NOT) & \texttt{!true} & \texttt{false} \\
\hline
\end{tabular}
\end{center}

\textbf{Beispiele:}
\begin{lstlisting}
int alter = 20;
boolean hatFuehrerschein = true;

// UND: Beide Bedingungen muessen wahr sein
if (alter >= 18 && hatFuehrerschein) {
    System.out.println("Darf Auto fahren");
}

// ODER: Mindestens eine Bedingung muss wahr sein
int note = 1;
if (note == 1 || note == 2) {
    System.out.println("Gute Note!");
}

// NICHT: Kehrt Wahrheitswert um
boolean gefunden = false;
if (!gefunden) {
    System.out.println("Noch nicht gefunden");
}

// Array-Grenzen pruefen
int[] array = {1, 2, 3};
int i = 1;
if (i >= 0 && i < array.length) {
    System.out.println(array[i]);
}

// Zeichen ist Buchstabe oder Ziffer
char c = 'A';
if ((c >= 'a' && c <= 'z') || (c >= 'A' && c <= 'Z')) {
    System.out.println("Ist ein Buchstabe");
}
\end{lstlisting}

\textbf{Kurzschluss-Auswertung:}
\begin{lstlisting}
// Bei && wird rechts nicht mehr ausgewertet, wenn links schon false ist
int[] zahlen = {1, 2, 3};
int index = 5;
if (index < zahlen.length && zahlen[index] > 0) {
    // Kein Fehler! zahlen[index] wird nicht ausgefuehrt, da index < zahlen.length bereits false ist
}
\end{lstlisting}

% ============================================================
\newpage


\section*{Zuweisungsoperatoren}
\addcontentsline{toc}{section}{Zuweisungsoperatoren}

\subsection*{Beschreibung:}\\
Weisen einer Variable einen Wert zu oder ändern ihn.

\begin{center}
\begin{tabular}{|l|l|l|l|}
\hline
\textbf{Operator} & \textbf{Bedeutung} & \textbf{Beispiel} & \textbf{Entspricht} \\
\hline
\texttt{=} & Zuweisung & \texttt{a = 5} & - \\
\hline
\texttt{+=} & Addition und Zuweisung & \texttt{a += 3} & \texttt{a = a + 3} \\
\hline
\texttt{-=} & Subtraktion und Zuweisung & \texttt{a -= 3} & \texttt{a = a - 3} \\
\hline
\texttt{*=} & Multiplikation und Zuweisung & \texttt{a *= 3} & \texttt{a = a * 3} \\
\hline
\texttt{/=} & Division und Zuweisung & \texttt{a /= 3} & \texttt{a = a / 3} \\
\hline
\texttt{\%=} & Modulo und Zuweisung & \texttt{a \%= 3} & \texttt{a = a \% 3} \\
\hline
\end{tabular}
\end{center}

\textbf{Beispiele:}
\begin{lstlisting}
int x = 10;
x += 5;   // x = 15 (entspricht: x = x + 5)
x -= 3;   // x = 12 (entspricht: x = x - 3)
x *= 2;   // x = 24 (entspricht: x = x * 2)
x /= 4;   // x = 6  (entspricht: x = x / 4)
x %= 4;   // x = 2  (entspricht: x = x % 4)

// Summe aufaddieren
int summe = 0;
int[] zahlen = {5, 10, 15};
for (int i = 0; i < zahlen.length; i++) {
    summe += zahlen[i];  // Kuerzer als: summe = summe + zahlen[i]
}

// Zeichen zaehlen
String text = "Hallo";
int anzahl = 0;
for (int i = 0; i < text.length(); i++) {
    if (text.charAt(i) == 'l') {
        anzahl++;  // Entspricht: anzahl = anzahl + 1
    }
}
\end{lstlisting}

% ============================================================
\newpage


\section*{Inkrement- und Dekrement-Operatoren}
\addcontentsline{toc}{section}{Inkrement- und Dekrement-Operatoren}

\subsection*{Beschreibung:}\\
Erhöhen oder verringern einen Wert um 1.

\begin{center}
\begin{tabular}{|l|l|l|l|}
\hline
\textbf{Operator} & \textbf{Bedeutung} & \textbf{Beispiel} & \textbf{Ergebnis} \\
\hline
\texttt{++} & Inkrement (erhöhen um 1) & \texttt{i++} & \texttt{i = i + 1} \\
\hline
\texttt{--} & Dekrement (verringern um 1) & \texttt{i--} & \texttt{i = i - 1} \\
\hline
\end{tabular}
\end{center}

\textbf{Beispiele:}
\begin{lstlisting}
int i = 5;
i++;        // i = 6 (entspricht: i = i + 1)
i--;        // i = 5 (entspricht: i = i - 1)

// In Schleifen
for (int j = 0; j < 5; j++) {  // j wird nach jedem Durchlauf um 1 erhoeht
    System.out.println(j);
}

// Countdown
int countdown = 10;
while (countdown > 0) {
    System.out.println(countdown);
    countdown--;
}

// Zaehler erhoehen
int anzahl = 0;
String text = "Hallo";
for (int k = 0; k < text.length(); k++) {
    if (text.charAt(k) == 'l') {
        anzahl++;
    }
}
\end{lstlisting}

\textbf{Unterschied: Prä- vs. Post-Inkrement:}
\begin{lstlisting}
int a = 5;
int b = a++;  // b = 5, dann a = 6 (Post-Inkrement: erst zuweisen, dann erhoehen)

int c = 5;
int d = ++c;  // c = 6, dann d = 6 (Prae-Inkrement: erst erhoehen, dann zuweisen)
\end{lstlisting}

% ============================================================
\newpage


\section*{String-Konkatenation}
\addcontentsline{toc}{section}{String-Konkatenation}

\subsection*{Beschreibung:}\\
Der \texttt{+} Operator verbindet Strings.

\textbf{Beispiele:}
\begin{lstlisting}
String vorname = "Anna";
String nachname = "Schmidt";
String vollname = vorname + " " + nachname;  // "Anna Schmidt"

// Mit Zahlen
int alter = 25;
String text = "Ich bin " + alter + " Jahre alt";  // "Ich bin 25 Jahre alt"

// In Schleifen
String wort = "Hallo";
String umgekehrt = "";
for (int i = wort.length() - 1; i >= 0; i--) {
    umgekehrt = umgekehrt + wort.charAt(i);
}
// umgekehrt = "ollaH"

// Array-Elemente verbinden
int[] zahlen = {1, 2, 3, 4, 5};
String ausgabe = "";
for (int i = 0; i < zahlen.length; i++) {
    ausgabe += zahlen[i];
    if (i < zahlen.length - 1) {
        ausgabe += ", ";
    }
}
// ausgabe = "1, 2, 3, 4, 5"
\end{lstlisting}

% ============================================================
\newpage


\section*{Operatorrangfolge (Priorität)}
\addcontentsline{toc}{section}{Operatorrangfolge (Priorität)}

\subsection*{Beschreibung:}\\
Operatoren werden in einer bestimmten Reihenfolge ausgewertet.

\textbf{Von höchster zu niedrigster Priorität:}

\begin{enumerate}
	\item \textbf{Klammern}: \texttt{( )}
	\item \textbf{Inkrement/Dekrement}: \texttt{++}, \texttt{--}
	\item \textbf{Negation}: \texttt{!}
	\item \textbf{Multiplikation, Division, Modulo}: \texttt{*}, \texttt{/}, \texttt{\%}
	\item \textbf{Addition, Subtraktion}: \texttt{+}, \texttt{-}
	\item \textbf{Vergleichsoperatoren}: \texttt{<}, \texttt{>}, \texttt{<=}, \texttt{>=}
	\item \textbf{Gleichheit}: \texttt{==}, \texttt{!=}
	\item \textbf{Logisches UND}: \texttt{\&\&}
	\item \textbf{Logisches ODER}: \texttt{||}
	\item \textbf{Zuweisung}: \texttt{=}, \texttt{+=}, \texttt{-=}, etc.
\end{enumerate}

\textbf{Beispiele:}
\begin{lstlisting}
int ergebnis = 5 + 3 * 2;      // ergebnis = 11 (nicht 16!)
// * wird vor + ausgefuehrt: 5 + (3 * 2) = 5 + 6 = 11

int x = (5 + 3) * 2;            // x = 16
// Klammern aendern die Reihenfolge: (5 + 3) * 2 = 8 * 2 = 16

boolean b = 5 > 3 && 2 < 4;     // b = true
// Erst Vergleiche, dann &&: (5 > 3) && (2 < 4) = true && true = true

int y = 10 + 5 * 2 - 3;         // y = 17
// Reihenfolge: 5 * 2 = 10, dann 10 + 10 = 20, dann 20 - 3 = 17

// Bei Unsicherheit: Klammern verwenden!
int z = ((10 + 5) * 2) - 3;     // z = 27 (macht Absicht deutlich)
\end{lstlisting}

\begin{achtung}
	\begin{itemize}
		\item Im Zweifelsfall immer Klammern setzen, um die Absicht klar zu machen!
		\item Multiplikation und Division haben höhere Priorität als Addition und Subtraktion
		\item Vergleichsoperatoren werden vor logischen Operatoren ausgewertet
		\item Klammern haben immer höchste Priorität
	\end{itemize}
\end{achtung}
